\chapter*{使用说明}\addcontentsline{toc}{chapter}{使用说明}
\section*{三条复习路线}
\begin{itemize}
\item \textbf{考研主线：}依次阅读第1--6章；先掌握定义、判定条件与方法卡，再独立重做基础例题和考研风格例题。章末每层前三题构成检测。
\item \textbf{竞赛路线：}在核心熟练后完成各章竞赛风格例题，并进入第7、8章。重点练习核像维数、秩估计、矩阵多项式与二次型不等式的证明。
\item \textbf{查缺路线：}用附录公式速查定位公式，以索引跳到定义与典型例题；随后回到章末习题，避免只看答案而没有独立演算。
\end{itemize}
\section*{阅读盒子与难度}
青绿色盒子放定义和方法，深蓝色盒子放定理和例题，金色盒子放考试提示与拓展，赭红色盒子放易错警示。颜色仅辅助区分，黑白打印仍可由标题识别。
$\star$表示直接应用定义或基本运算；$\star\star$表示需要选择方法或分类；$\star\star\star$表示结构证明与综合迁移。星级与考研、竞赛标签相互独立。

每章四节，每节四例。标注“方法一、方法二”的例题应至少亲手完成其中一种，再解释另一种何时更合适。章末九题依次为基础、考研风格和竞赛拓展各三题。题目后的链接通往答案，答案末尾可返回题目。
\section*{训练闭环}
第一次只读题目与思路，写出条件和目标；第二次完成推导并检查维数、数域、可逆性；第三次比较方法点评并处理变式。错题记录应写成可执行规则，例如“求原列空间的基时，要返回原矩阵的主元列”，而不是“计算要细心”。

建议把一天的练习分成概念复述、限时计算和一道证明三部分。时间可按个人基础调整，不把固定进度当作掌握程度。最后复习时，优先检查含参数的分支、可对角化的重根判断、相似与合同的区别、正定与半正定判据。
\section*{如何读图与使用巧解}
先看图中的向量、直线或等值线分别对应哪个代数对象，再看操作前后哪一个量保持不变。图旁的“理解札记”解释结构，“巧解与数形结合”把结构转成可操作步骤。读完应能把图重新翻译成方程或维数关系，而不只是记住图形外观。

在二维图上成立的直觉，推广到高维时应由定理保证。例如两直线相交的图像能提示秩判据，但不能代替一般相容性的证明；椭圆能帮助理解正定型，但非对称矩阵不能直接套用实对称谱定理。每次使用捷径都检查盒子中写出的条件。

可从图解目录直接跳转到17幅图，再通过相邻例题完成计算训练。章末速查表适合临考复盘，答案附录适合检查参数、符号、自由度与等号条件。
